\begin{center}
\hypertarget{native-036-article018}{}\pdfbookmark[2]{18.}{article-018}18.
\end{center}
In PARABOLA anomalia excentrica, anomalia media atque motus medius fierent $=0$; hic igitur istae notiones comparationi motus cum tempore inservire nequeunt. Attamen in parabola angulo auxiliari ad integrandum $r r\,\mathrm{d}v$ omnino opus non habemus; fit enim
\[
\begin{aligned}
r r\,\mathrm{d}v&=\frac{p p\,\mathrm{d}v}{4\cos\frac{1}{2}v^4}=\frac{p p\,\mathrm{d}\operatorname{tang}\frac{1}{2}v}{2\cos\frac{1}{2}v^2}\\
&=\frac{1}{2}p p(1+\operatorname{tang}\frac{1}{2}v^2)\,\mathrm{d}\operatorname{tang}\frac{1}{2}v
\end{aligned}
\]
adeoque $\int r r\,\mathrm{d}v=\frac{1}{2}p p(\operatorname{tang}\frac{1}{2}v+\frac{1}{3}\operatorname{tang}\frac{1}{2}v^3)+\mathrm{Const.}$ Si tempus a transitu per perihelium incipere supponitur, Constans fit $=0$; habetur itaque
\[
\operatorname{tang}\frac{1}{2}v+\frac{1}{3}\operatorname{tang}\frac{1}{2}v^3=\frac{2t k\BareRoot{(1+\mu)}}{p^{\frac{3}{2}}},
\]
\SourceLeaf{037}per quam formulam $t$ ex $v$, atque $v$ ex $t$ derivare licet, simulac $p$ et $\mu$ sunt cognitae. Pro $p$ inter elementa parabolica radius vector in perihelio qui est $\frac{1}{2}p$ exhiberi, massaque $\mu$ omnino negligi solet. Vix certe unquam possibile erit, massam corporis talis, cuius orbita tamquam parabola computatur, determinare, reveraque omnes cometae per optimas recentissimasque observationes densitatem atque massam tam exiguam habere videntur, ut haec insensibilis censeri tutoque negligi possit.
